\selectlanguage{ngerman}%
\frenchspacing%
\begin{ledgroupsized}[r]{120mm}%
\footnotesize%
\pstart%
\noindent\textbf{Überlieferung:}%
\pend%
\end{ledgroupsized}%
\begin{ledgroupsized}[r]{114mm}%
\footnotesize%
\pstart%
\parindent -6mm%
\makebox[6mm][l]{\textit{L}}%
Aufzeichnung:
LH XXXV 10,\;14 Bl.\;4.
Ein Blatt 8\textsuperscript{o},
Ränder zum Teil gerissen;
Fragment eines \href{\Scan{LH_35_10_14_004wz}}{Wasserzeichens}:
niedersächsisches Papier. %
Zwei Drittel Seite auf Bl.\;\href{\Scan{LH_35_10_14_004r}}{4\,r\textsuperscript{o}},
unteres Drittel unbeschrieben;
Bl.\;\href{\Scan{LH_35_10_14_004v}}{4\,v\textsuperscript{o}} überliefert nur einen abbrechenden und gestrichenen Satz von Leibnizens Hand:
\textit{Si duo corpora inter se cohaereant, vis respectiva ipsorum et vim respectivam inter se co}.
Am oberen Rand von Bl.\;\href{\Scan{LH_35_10_14_004r}}{4\,r\textsuperscript{o}} liegt, von Leibnizens Hand und rot unterstrichen, die Feststellung \textit{Hoc recte} vor;
daneben von unklarer Hand und in roter Tinte der Vermerk \textit{:\:IV}
(vgl. die Träger der benachbarten Signaturen LH~XXXV~10,\;14 Bl.\;2 bis~13,
u.\,a. 
N.\;\ref{60117},
\ref{60118},
\ref{60115},
\ref{60112},
\ref{60116},
\ref{60113},
\ref{60114},
\ref{60110}
u.\;\ref{41824}
im vorliegenden Band).
\pend%
\end{ledgroupsized}%
\vspace{5mm}%
\begin{ledgroup}%
\footnotesize%
\pstart%
\noindent%
\textbf{Datierungsgründe:}
Die Datierung der Aufzeichnung N.\;\ref{60111} beruht auf Eigenschaften des Textträgers sowie auf inhaltlichen Betrachtungen.
Siehe die ausführliche Begründung auf S.\;\pageref{60111_dat}.%
\pend%
\end{ledgroup}%
\selectlanguage{latin}%
\frenchspacing%
\vspace{8mm}%
\normalsize%
\pstart%
\noindent%
\bracks{%
\href{\Scan{LH_35_10_14_004r}}%
{4\,r\textsuperscript{o}}%
}
\edtext{%
Quam vim respectivam%
\indexS{vis respectiva conversa in circulationem}
duo corpora inter se habitura essent,
si%
}{%
\lemma{\bracks{%
\href{\Scan{LH_35_10_14_004r}}{4\,r\textsuperscript{o}}%
}}%
\Bfootnote{%
\textit{(1)}~\angles{Si duo} %
\textit{(2)}~Quam vim %
\bracks{...} essent, si%
~\textit{L}%
}}
non firmiter cohaererent,%
\indexS{corpora non cohaerentia}
eam in casu firmae cohaesionis convertent%
\indexS{cohaesio firma}%
\indexS{vis respectiva conversa in circulationem}
in circulationem circa commune centrum gravitatis.%
\indexS{circulatio circa centrum gravitatis}%
\indexS{centrum gravitatis commune}
\pend%
\begin{edfigure}[0.50em][1.00em]%
\hfill%
\includegraphics[scale=0.87]%
{gesamttex/edit_VIII-4/images/LH_35_10_14_004_d1_004r.pdf}%
\hfill%
\includegraphics[scale=0.80]%
{gesamttex/edit_VIII-4/images/LH_35_10_14_004_d2_004r.pdf}%
\hfill%
\edcaption[0.75em]{%
\lbrack\textit{Fig.\:1}\rbrack%
\hspace*{46mm}%
\lbrack\textit{Fig.\:2}\rbrack%
}
\end{edfigure}%
\pstart%
Finge \textit{A} in \textit{B} incurrere,%
\indexS{corpora incurrentia sibi}
si jam \textit{B} non cohaereat ipsi \textit{C},%
\indexS{corpora non cohaerentia}
patet
quae inter \textit{B} et \textit{C}
futura sit respectiva vis%
\indexS{vis respectiva conversa in circulationem}
seu a se invicem recedendi conatus.%
\indexS{conatus recedendi a se invicem}
Sed cum ob firmitatem%
\indexS{firmitas corporis}
nequeant a se invicem recedere,
necesse est
ut vis respectiva%
\indexS{vis respectiva conversa in circulationem}
debita inter ipsa servetur, %
eam verti in circulationem%
\indexS{circulatio circa centrum gravitatis}
circa centrum gravitatis,%
\indexS{centrum gravitatis commune} %
ita enim eadem semper vi a se invicem recedere conabuntur.
Ipsum autem centrum commune%
\indexS{centrum gravitatis commune}
\ledpb%
eadem celeritate%
\indexS{celeritas centri gravitatis}
recedet ab \textit{A},
qua recessisset sine cohaesione.%
\indexS{cohaesio corporum}
Itaque nulla mutatio est,%
\indexS{mutatio virium}
nisi quod pars virium%
\indexS{pars virium translata}
transfertur in \textit{C},%
\indexS{pars virium translata}
ut eadem maneat vis respectiva%
\indexS{vis respectiva conversa in circulationem}
inter \textit{B} et \textit{C}
quae alias,
et inter \textit{A} et \textit{BC} %
quae alias.
Hinc patet
etiam vim centrifugam%
\indexS{vis centrifuga}
non designare vim absolutam,%
\indexS{vis absoluta}
sed solum vim respectivam,%
\indexS{vis respectiva conversa in circulationem}
adeoque nil mutare,
nec opus ejus
\edtext{%
determinatam Hypothesin%
}{%
\lemma{determinatam}%
\Bfootnote{%
\bars{ejus \textit{gestr.}} Hypothesin%
~\textit{L}%
}}
celeritatis%
\indexS{hypothesis celeritatis}
posse haberi%
\bracks{,}
\edlabel{LH_35_10_14_004_004r_Newt-1}%
\edtext{%
quod Neutonus minus observavit.%
\indexN{\textls{Newton} (Neutonus), Isaac 1643\textendash1727}%
}{%
\lemma{quod \bracks{...} observavit}%
\Cfootnote{%
Möglicherweise Anspielung auf
\textsc{I.\;Newton}, \textit{Principia},
lib.\;I, sect.\;II, prop.\;IV, theor.\;IV, scholium
(London 1687, S.\;43).%
\cite{00535}%
}}%
\edlabel{LH_35_10_14_004_004r_Newt-2}
Quantitas autem corporis \textit{B} definitur,%
\indexS{quantitas corporis}
\edtext{%
\edtext{%
si \textit{AE}%
}{%
\lemma{si}%
\Bfootnote{%
\bars{si \textit{streicht Hrsg.}} \textit{AE}%
~\textit{L}%
}}
normalis%
}{%
\lemma{si \textit{AE} normalis}%
\Cfootnote{%
Vgl. das Diagramm \bracks{\textit{Fig.\:2}}.%
}}
ad tangentem superficiei producatur%
\indexS{normalis ad tangentem superficiei}
et maxima corporis pars%
\indexS{centrum partis corporis maximae}%
\bracks{,}
cujus centrum cadit in \textit{EB} productam%
\bracks{,}
erit %
ea
quae centrice percutitur;%
\indexS{pars corporis percussa centrice}
reliqua vero non
nisi ob firmitatem.%
\indexS{firmitas corporis}
\pend%
\pstart%
Sed hic illud notandum,
fieri posse
\edtext{%
ut haec recta \textit{EB}
sit obliqua ad \textit{CB}
jungentem centra gravitatis.%
\indexS{recta jungens centra gravitatis}%
\indexS{centra gravitatis juncta}%
}{%
\lemma{\textit{Am Rand:}}%
\Afootnote{%
(fig.\:2)%
\indexS{figura}%
}}
Sit \textit{ED} parallela \textit{CB},
et angulus \textit{EDB} rectus,%
\indexS{angulus rectus}
\edtext{%
circulatio circa centrum commune%
\indexS{circulatio circa centrum gravitatis}%
\indexS{centrum gravitatis commune}
ipsorum \textit{C} et \textit{B}
tantum%
}{%
\lemma{circulatio}%
\Bfootnote{%
\textit{(1)}~tantum %
\textit{(2)}~circa centrum %
\bracks{...} et \textit{B} tantum%
~\textit{L}%
}}
imprimitur conatu \textit{DB},%
\indexS{conatus circulationis impressae}
non vero conatu \textit{ED};%
\indexS{conatus circulationis impressae}
quo perinde est
ac si incurrens%
\indexS{corpus incurrens}
in
\edtext{%
ipsa\bracks{m} recta\bracks{m}%
}{%
\lemma{ipsa recta}%
\Bfootnote{%
\textit{L\;ändert Hrsg.}%
}}
\textit{CB} moveretur,
et divellere conaretur \textit{B} et \textit{C},
adeoque ob eorum firmitatem%
\indexS{firmitas corporis}
hinc toti impetum conatu \textit{ED} imprimeret.%
\indexS{impetus impressus}%
\indexS{conatus impetus impressi}
\pend%
\begin{edfigure}[2.0em][1.25em]%
\includegraphics[scale=0.87]%
{gesamttex/edit_VIII-4/images/LH_35_10_14_004_d3_004r.pdf}%
\edcaption{%
\lbrack\textit{Fig.\:3}\rbrack%
\edtext{}{%
\lemma{\lbrack\textit{Fig.\:3}\rbrack}%
\Cfootnote[nonum]{%
\textit{L} überliefert (Bl.\;\href{\Scan{LH_35_10_14_004r}}{4\,r\textsuperscript{o}}) einen gestrichenen, hier nicht wiedergegebenen Entwurf zum Diagramm.%
}}%
}%
\end{edfigure}%
\pstart%
Poterit etiam hac methodo determinari%
\indexS{methodus determinandi}
quid fiat
si plura simul incurrant in idem%
\indexS{incursus plurium in idem}
eccentrice diverso modo,%
\indexS{incursus eccentricus}
inde diversae fiunt simul divisiones.%
\indexS{divisio corporis ob incursum}
Dividamus rem in tot partes,
quot opus est,
ut quaevis divisio%
\indexS{divisio corporis ob incursum}
\edtext{%
succedat.
Ut%
}{%
\lemma{succedat.}%
\Bfootnote{%
\textit{(1)}~Et %
\textit{(2)}~Ut%
~\textit{L}%
}}
si fig.\;\angles{3}%
\indexS{figura}
unus incursus dividat totum \textit{AC}%
\indexS{incursus dividens corpus}
in partes \textit{AB}, \textit{BC},
alter in partes \textit{AD}, \textit{DC},
habemus divisionem in tres partes
\textit{AB}, \textit{BD}, \textit{DC},
ubi \textit{BD} habet duos impulsus communes.%
\indexS{impulsus communis}
\pend%
\newpage%
\pstart%
\noindent%
\centering%
*     *     *
\pend%
\vspace{2.5em}%
\selectlanguage{german}%
\frenchspacing%
\footnotesize%
\pstart%
\noindent%
\textls{Ausführliche Datierungsgründe:}%
\pagelabel{60111_dat}
Der Träger der titellosen und undatierten Aufzeichnung N.\;\ref{60111} weist das Fragment eines Wasserzeichens auf, das von der Papiermühle Relliehausen am Solling\indexO{Relliehausen am Solling} (Niedersachsen)\indexO{Niedersachsen} stammte. %
Papier mit diesem Wasserzeichen hat Leibniz nach heutigem Wissensstand in der frühen Hannoveraner\indexO{Hannover} Zeit verwendet, d.\,h. vor seiner Reise nach Süddeutschland,\indexO{Süddeutschland} Österreich\indexO{Österreich} und Italien\indexO{Italien} zwischen November 1687 und der ersten Hälfte Juni 1690.
Ausschlaggebend für die Datierung der Aufzeichnung N.\;\ref{60111} aber ist vielmehr der dort anzutreffende Verweis auf Newtons\indexN{\textls{Newton} (Neutonus), Isaac 1643\textendash1727} Betrachtung der zentrifugalen Kraft (\textit{vis centrifuga}) in den 1687 veröffentlichten \textit{Principia}\cite{00535} (siehe oben, \refpassage{LH_35_10_14_004_004r_Newt-1}{LH_35_10_14_004_004r_Newt-2}).
Der Verweis zeigt, dass N.\;\ref{60111} nicht vor Leibnizens eigenständiger Lektüre von Newtons\indexN{\textls{Newton} (Neutonus), Isaac 1643\textendash1727} Werk entstanden sein kann.
Die von Christoph Pfautz,\indexN{\textls{Pfautz} (Pfauzius), Christoph 1645\textendash1711} veröffentlichte Rezension\cite{01356} der \textit{Principia}\cite{00535} (\textit{AE},\cite{01023} Juni 1680, S.\;303--315), die Leibniz gelesen und exzerpiert hat (vgl. LH~XXXV~10,\;7 Bl.\;13--14), kann nicht die Quelle für Leibnizens Verweis sein, weil dort die zentrifugale Kraft unerwähnt bleibt.
\pend%
\pstart%
Leibnizens erste bezeugte Lektüre von Newtons\indexN{\textls{Newton} (Neutonus), Isaac 1643\textendash1727} \textit{Principia}\cite{00535} fand -- wie \textsc{Bertoloni Meli} 1993\cite{01357} (S.\;97) nachgewiesen hat -- im Herbst 1688 statt, als er sich, auf dem Weg nach Italien,\indexO{Italien} in Wien\indexO{Wien} aufhielt.
Der Herbst 1688 ist folglich als frühester Entstehungszeitpunkt der Aufzeichnung N.\;\ref{60111} anzusehen.
Bei deren Abfassung verwendete Leibniz offenbar älteres Papier aus Niedersachsen.\indexO{Niedersachsen}
Dies macht eine Entstehung der Aufzeichnung N.\;\ref{60111} nach Mitte Juni 1690, als Leibniz nach Hannover\indexO{Hannover} zurückkam, wahrscheinlicher;
eine frühere Abfassung des Textes zwischen Wien\indexO{Wien} und Italien\indexO{Italien} ist jedoch nicht auszuschließen (Leibniz kann im November 1687 Papier aus Hannover\indexO{Hannover} mitegebracht und später auf der Reise verwendet haben).
\pend%
\pstart%
Für die Datierung der Aufzeichnung N.\;\ref{60111} lässt sich zum Redaktionsschluss kein eindeutiger Terminus ante quem angeben.
\pend%
\selectlanguage{latin}%
\frenchspacing%
\normalsize%
